\documentclass[openany]{ctexbook}
\usepackage{amsmath,amssymb}
\usepackage[hidelinks]{hyperref}
\newcommand{\note}[1]{\footnote{译注：#1}}
\title{代数几何基础第五卷预备笔记\\10.2及其后编排说明}
\author{据用户上传英译编辑稿译出}
\date{2026年11月15日}
\begin{document}
\maketitle
\chapter*{范围与校订状态}
本批连续译出底本第39页10.2“By a polarity”起至第40页命题10.3前的编排批注，含原脚注56。不包含命题10.3的陈述和证明。正文保留草稿的论断与写作计划，未将其升级为已独立验证的数学结果；对于基域、有理性及正则性约定的疑点作译注。此前8.8缺文、9.4待陈述、9.5残文及10.1疑点均未解决。本批为独立增补，未合入原合订本。
\chapter*{10.2 线性束与爆破}
借助极性对应，$P^v$ 中的一条直线 $L$ 对应于 $P$ 中余维为2的线性子簇 $L^0$（$S$ 任意）。置
\[
T=X\times_P L^0.
\]
\note{底本此处印作“$T=x_P L^0$”，与随后零点概形的描述相结合，可明确读为沿 $X\to P$ 的逆像，即上述纤维积；这里补全纤维积符号并统一 $X$ 的大小写。}
还可以如下描述 $T$：$L$ 由 $E^v$ 的一个秩为2的局部自由商模定义；等价地，由 $E$ 的一个处处秩为2、局部为直和因子的子模 $F$ 定义。考虑复合同态
\[
F_X\longrightarrow E_X\longrightarrow\mathcal O_X(1),
\]
则 $T$ 正是这个复合同态的零点概形；换言之，它由理想 $J$ 定义，其中 $J$ 是相应同态（通过张量乘以 $\mathcal O_X(-1)$ 得到）
\[
F_X(-1)\longrightarrow\mathcal O_X
\]
的像。

假设这个同态是正则的，意思是：在局部写出 $F_X(-1)$ 的一个全序基，其在 $\mathcal O_X$ 中的像构成一个 $\mathcal O_X$-正则序列。这个条件与所选的基无关；也可以内在地表述为：
\[
F_X(-1)\otimes_{\mathcal O_X}(\mathcal O_X/J)
\longrightarrow J/J^2
\]
是同构，并且 $V(J)=T\to X$ 是正则浸入。\note{这里忠实保留原稿的正则性表述；在任意预概形上，所用“正则序列”的局部约定、是否仅在零点概形附近要求该条件，以及更换所选局部基后的不变性与上述内在表述的确切等价范围，应与作者有关正则浸入的一节统一核查。本批未独立证明这一等价。}
〔原稿附记：应当在某处阐明一般情形，即预概形 $X$ 上有一个同态 $G\to\mathcal O_X$，其中 $G$ 局部自由；例如可放在关于正则浸入的一节。〕于是有：

\noindent\textbf{定理10.2.}
在上述假设下，带有典范投影 $Y_L\to X$ 的线性束 $Y_L$，以唯一的方式在 $X$ 上同构于预概形 $X$ 沿中心 $T$ 的爆破。

为了理解这个定理的意义，宜在本节或本条开头指出：若 $S=\operatorname{Spec}(k)$，则对 $P^v$ 中一条“充分一般”的直线 $L$，上述正则性条件成立（参见第12节所指的补述，即关于5.3的补述）。\note{这是原稿对后文的预告，不应理解为无须约定定义域便可在任意有限域上选取一条所需的 $k$-有理直线；其精确范围仍须结合5.3与第12节核对。}
在下文构造本条开头所指出或预期的“优良”线性束时，我们可以要求所述束满足这一条件（它与下述条件属于同一类型，但并不相同：对每个 $\xi\in L$，要求 $\phi_\xi$ 是 $X_{k(\xi)}$-正则的）。应把所讨论的条件纳入拟作的汇总性陈述。

另一方面，实际上10.2只用于9.5的情形，因此，最好在10.2之后才用束的语言重新表述9.5，以便在该陈述中同时纳入束与一个爆破的同构（即给出便于适当引用的情形描述）。这样，对于无限域 $k$ 上每一个维数 $\ge2$ 的光滑、几何连通的射影概形 $X$，我们便得到一种方法，找出一个非空的光滑闭子概形，它在其每一点的余维均为2，使得爆破后的概形具有到 $\mathbf P^1$ 的一个纤维化：所有纤维均几何整，并且除至多有限个纤维外，其余纤维都光滑；那些例外纤维至多有一个几何奇异点，而且这样的点在 $k$ 上有理，并且在几何意义下是普通奇点。\note{底本确实写“rational over $k$”，此处没有改成参数点的剩余域。纤维的定义域与奇异点的有理性必须区分；若参数点不在 $k$ 上有理，不能径直把相应剩余域上的有理性写成 $k$-有理性。由于这里所引用的9.5在现有底本中残缺，本批将整段视为原稿的预期应用并保留该处待核，不声称已经补齐9.5或证明上述存在性。}

这说明，深入研究这种带奇异纤维的纤维化是重要的（这种研究在目前才刚刚开始）：它可在一定程度上把维数为 $d$ 的光滑射影簇的研究，归约为对维数为 $d-1$、容许有普通奇点的射影簇的研究（这些簇组成依赖一个参数的族）。\note{底本在此句有未闭合括号及语法不顺的“or”；这里依可辨的语义整理括号与语序，没有加入新的数学结论。}

陈述10.2或多或少是下面命题的直接推论；后者完全独立于超平面截面这一主题，无疑更适合放在另设的一段“正则浸入”中。\footnote{原脚注56：“请询问A.G.”}

〔底本编排批注〕命题10.3被划去了。请询问A.G.，这是否是他的本意。
\end{document}
